\documentclass[11pt,a4paper]{article}
\usepackage[T1]{fontenc}
\usepackage{lmodern,amsmath,amssymb,amsthm,mathtools,booktabs,microtype}
\usepackage[margin=26mm]{geometry}
\usepackage[colorlinks=true,linkcolor=blue,citecolor=blue,urlcolor=blue]{hyperref}
\newtheorem{theorem}{Theorem}[section]
\newtheorem{proposition}[theorem]{Proposition}
\newtheorem{lemma}[theorem]{Lemma}
\newtheorem{corollary}[theorem]{Corollary}
\theoremstyle{definition}
\newtheorem{definition}[theorem]{Definition}
\newtheorem{remark}[theorem]{Remark}
\newcommand{\Q}{\mathbb Q}\newcommand{\Z}{\mathbb Z}
\newcommand{\A}{\mathbb A}\newcommand{\PP}{\mathbb P}
\newcommand{\Gm}{\mathbb G_m}\newcommand{\Ok}{\mathcal O_k}
\DeclareMathOperator{\Dom}{Dom}
\setlength{\emergencystretch}{2em}
\title{Finite Correspondences over Abelian Varieties\newline and Bounded Denominators}
\author{}\date{9 November 2026}
\begin{document}\maketitle
\begin{abstract}
Let a variety over a number field map quasi-finitely to an abelian
variety. We prove that the values of any rational function on its
rational points meet each finitely generated multiplicative subgroup
in a finite set. The same conclusion holds whenever the function
takes only finitely many values on each geometric fiber of a morphism
to an abelian variety. This consequence of semiabelian Mordell--Lang
extends the rational-map obstruction to finite algebraic
correspondences. It implies omission of almost every term of every
fixed affine geometric progression, and weak approximation for the
complement. For generically finite morphisms, sufficiently late
represented terms must come from the exceptional locus. An explicit
smooth surface shows why this exception is necessary. We also exclude
bounded-denominator sets containing all integers when the witness
fibers are soluble unirational varieties, or form a smooth projective
family of genus-one curves over an abelian base. The latter argument
uses a divisor-class map to the relative Jacobian, Siegel's theorem, and
Merel's uniform torsion theorem. These conclusions identify restrictions
on constructions of Diophantine sets; they do not prove undefinability
of a ring of integers in its field of fractions.
\end{abstract}

\section{Definitions and the main statement}

Throughout, $k$ is a number field and $\overline k$ is an algebraic
closure. Varieties are reduced separated schemes of finite type.
A rational function is evaluated on a specified open domain on which
it is a morphism to $\A^1$. Restricting that domain is permitted.
For a morphism $\pi:V\to A$ and $a\in A(\overline k)$, write
$V_a$ for its geometric fiber. A morphism is \emph{quasi-finite}
if it is of finite type and all its geometric fibers are finite.

\begin{definition}
A subset $T\subset k$ has \emph{uniformly bounded denominators}
if there is a positive integer $M$ such that $MT\subseteq\Ok$.
It has the \emph{weak-approximation property} if its diagonal image
is dense in $\prod_{v\in S}k_v$ for every finite set $S$ of places
of $k$.
\end{definition}

\begin{theorem}\label{thm:main}
Let $A/k$ be an abelian variety, let $\pi:V\to A$ be a morphism,
and let $h:U\to\A^1_k$ be the restriction of a rational function
to an open subset $U\subseteq V$. Suppose that
\begin{equation}\label{eq:finite-values}
       h\bigl(U(\overline k)\cap V_a(\overline k)\bigr)
       \quad\hbox{is finite for every }a\in A(\overline k).
\end{equation}
Then, for every finitely generated subgroup $\Lambda\subset k^*$,
the set
\[
                         h(U(k))\cap\Lambda
\]
is finite. In particular, the conclusion holds if $\pi$ is
quasi-finite.
\end{theorem}

The condition is about geometric fibers. Finiteness of the sets
$U(k)\cap V_a(k)$ alone is not the stated hypothesis. Finite covers
of arbitrary degree satisfy the hypothesis; no rational section or
point-stability assertion is required.

For rational maps on abelian varieties, Corvaja--Zannier prove that
finite unions of their rational value sets omit infinitely many
integers and have a complement satisfying weak approximation
\cite[Theorems 1.1 and 1.10]{CZ}. Theorem~\ref{thm:main} permits
outputs that are algebraic over the abelian coordinates, provided
their geometric choices remain finite at every relevant point.

\section{The constructible-image argument}

We first record the arithmetic theorem in the form needed here.
A semiabelian variety is an extension of an abelian variety by a
torus.

\begin{theorem}[Semiabelian Mordell--Lang]\label{thm:ML}
Let $G$ be a semiabelian variety over a number field and let
$\Gamma\subset G(\overline k)$ be finitely generated. The Zariski
closure of any subset of $\Gamma$ is a finite union of translates
of connected algebraic subgroups of $G$.
\end{theorem}

This is a standard closure form of the semiabelian theorem
\cite{VojtaI,VojtaSurvey}. Indeed, if $Z$ is the closure of the
given subset, the usual statement applies to $Z\cap\Gamma$;
that intersection is dense in $Z$. Taking closures of its finitely
many subgroup cosets gives the assertion. One may equivalently use
the integral-point form in \cite[Theorem 2.2]{VojtaSurvey}, after
choosing one ring of $S$-integers containing the finitely many group
generators and their inverses in a group model. Splitting a
disconnected subgroup into its finitely many connected components
gives the stated form. All subsequent subgroup arguments may be
made over $\overline k$.

\begin{lemma}\label{lem:subgroups}
Every connected algebraic subgroup of $A\times\Gm$, where $A$
is an abelian variety over an algebraically closed field of
characteristic zero, is of one of the forms
\[
                         B\times\{1\},\qquad B\times\Gm,
\]
where $B$ is an abelian subvariety of $A$.
\end{lemma}
\begin{proof}
Let $H$ be such a subgroup, and let $B$ be its image in $A$.
The image of a homomorphism of algebraic groups is closed, so $B$
is an abelian subvariety. The kernel of $H\to B$ is a subgroup
of $\Gm$, hence is either finite or all of $\Gm$.

In the latter case each nonempty fiber contains every multiplicative
coordinate, giving $H=B\times\Gm$. In the former case the
surjective homomorphism with finite kernel is finite. Thus $H$ is
proper because $B$ is proper. Every morphism from the connected
proper variety $H$ to $\Gm$ is constant. Its value at the identity
is $1$, so $H=B\times\{1\}$.
\end{proof}

\begin{lemma}\label{lem:constructible}
Let $W\subset A\times\Gm$ be constructible, with finite geometric
fibers under the first projection. If $\Lambda\subset k^*$ is
finitely generated and
\[
                  \Sigma\subset W(k)\cap(A(k)\times\Lambda),
\]
then the second coordinates of the points of $\Sigma$ form a finite
set.
\end{lemma}
\begin{proof}
The Mordell--Weil theorem makes $A(k)\times\Lambda$ finitely
generated. Apply Theorem~\ref{thm:ML} to the closure
$Z=\overline\Sigma$. It is essential here to close this actual
arithmetic subset, rather than all of $W$.

After splitting components, every irreducible component of $Z$
is a translate of one of the groups in Lemma~\ref{lem:subgroups}.
Suppose that one component is $(a+B)\times\Gm$. The subset
$\Sigma$ is dense in this component after the other irreducible
components have been removed. Hence its constructible intersection
with $W$ is dense and contains a nonempty open subset of
$(a+B)\times\Gm$. Choose a geometric point in that open subset.
Its intersection with the vertical fiber through the point is a
nonempty open subset of $\Gm$, and is therefore infinite. This
contradicts the finite-fiber hypothesis on $W$.

Every component of $Z$ consequently has a constant second
coordinate. There are finitely many components, proving the
assertion. Translates by torsion points and disconnected groups
introduce only finitely many such components and cause no exception.
\end{proof}

\begin{proof}[Proof of Theorem~\ref{thm:main}]
Restrict $U$ to $h^{-1}(\Gm)$ and take its image $W$ under
$(\pi,h)$. Chevalley's theorem makes $W$ constructible
\cite[Tag 054K]{Stacks}. Hypothesis~\eqref{eq:finite-values}
says exactly that its geometric fibers over $A$ are finite.
Apply Lemma~\ref{lem:constructible} to the image points coming
from $U(k)$ with $h$-value in $\Lambda$. The lemma allows such
a subset; it does not assume that every rational point of the
constructible image has a rational lift to $U$.
\end{proof}

\section{Arithmetic progressions and weak approximation}

\begin{corollary}\label{cor:affine-units}
Let $T$ be a finite union of value sets satisfying
Theorem~\ref{thm:main}, together with a finite subset of $k$.
For every $a,b\in k$ with $b\ne0$, and every finitely generated
$\Lambda\subset k^*$,
\[
                           T\cap(a+b\Lambda)
\]
is finite. In particular, for every integer $q\ge2$, all but
finitely many terms $a+bq^n$, $n\ge0$, are outside $T$.
\end{corollary}
\begin{proof}
Replace each output $h$ by $(h-a)/b$. This affine change preserves
finiteness of its values on each geometric fiber, so
Theorem~\ref{thm:main} applies. A finite union and a finite added
set preserve finiteness. Take $\Lambda$ generated by $q$ for the
last assertion.
\end{proof}

\begin{corollary}\label{cor:weak}
The complement $k\setminus T$ in Corollary~\ref{cor:affine-units}
satisfies weak approximation.
\end{corollary}
\begin{proof}
Fix neighborhoods at finitely many places. Weak approximation in
$k$ supplies a point $a\in k$ inside all of them. Choose a nonzero
$b\in k$ sufficiently divisible at the prescribed finite places,
so that adding any element of the same absolute values as $b$ at
those places keeps the point in the neighborhoods. Choose a rational
prime $q$ below none of these finite places.

The values $a+bq^{-n}$ are distinct. At each prescribed finite
place, $q$ is a unit, so the chosen divisibility of $b$ works for
every $n$. At every archimedean place, $bq^{-n}$ tends to zero.
Thus all sufficiently late terms lie in the prescribed product of
neighborhoods. Only finitely many belong to $T$, by
Corollary~\ref{cor:affine-units} applied to the group generated
by $q$. The remaining terms prove the assertion.
\end{proof}

For $a,b\in\Z$, $b\ne0$, the omitted values $a+b2^n$ lie in
the integer arithmetic progression $a+b\Z$. Thus none of these
constructions contains any complete integer arithmetic progression.
In particular it cannot supply a set $S$ with
\begin{equation}\label{eq:sandwich}
                         \Z\subseteq S,\qquad MS\subseteq\Ok
\end{equation}
for a fixed positive integer $M$. This obstruction concerns the
first inclusion already; it needs no denominator hypothesis.
The proof supplies no effective bound on the omitted exponents.

\section{Algebraic outputs and exceptional loci}

\begin{proposition}\label{prop:algebraic-output}
Let $Y$ be a locally closed subvariety of an abelian variety $A/k$,
let $c_0,\ldots,c_d$ be regular functions on $Y$, and put
\[
              C=\{a\in Y:c_0(a)=\cdots=c_d(a)=0\}.
\]
Suppose $\pi:V\to Y$ and a regular output $h:U\to\A^1$ satisfy
\[
                  \sum_{j=0}^{d}c_j(\pi(v))h(v)^j=0
                  \qquad(v\in U).
\]
For each fixed $a,b\in k$, $b\ne0$, and finitely generated
$\Lambda\subset k^*$, only finitely many values in $a+b\Lambda$
have a witness with $\pi(v)\notin C$. Every other represented
value in that set has all its witnesses above $C$.
If the relation is monic in $h$, its full value set misses all
but finitely many terms of every affine geometric progression.
\end{proposition}
\begin{proof}
Above $Y\setminus C$ the polynomial in the output is nonzero,
so it has finitely many geometric roots. This gives
\eqref{eq:finite-values} for the composite map to $A$, regardless
of how many further witnesses lie above each point. Apply
Corollary~\ref{cor:affine-units}. A value outside that finite
set cannot have any witness above $Y\setminus C$. In the monic
case $C$ is empty.
\end{proof}

Thus adding unrestricted auxiliary variables cannot repair a monic
algebraic output relation over abelian coordinates. For a relation
with varying leading coefficient, the common-zero locus of all
coefficients must be checked: there the relation imposes no
condition on the output.

\begin{proposition}\label{prop:exceptional}
Let $V$ be integral and let $\pi:V\to A$ be generically finite
onto its image. Let $V_{\mathrm{qf}}$ be its quasi-finite locus and
$E=V\setminus V_{\mathrm{qf}}$. Then $V_{\mathrm{qf}}$ is a dense
open subset. For a rational output $h$, all but finitely many
represented values in any fixed set $a+b\Lambda$ have all their
witnesses in $E$.
\end{proposition}
\begin{proof}
The quasi-finite locus is open, and it contains the generic point
when the extension of function fields is finite
\cite[Tag 02FW]{Stacks}. Restricting $h$ to its intersection with
that locus gives a value set satisfying Theorem~\ref{thm:main}.
Apply Corollary~\ref{cor:affine-units}.
\end{proof}

One can further stratify $E$ into finitely many locally closed
subvarieties and repeat the argument on each stratum whose map to
$A$ is quasi-finite. Such strata account for only finitely many
values in a fixed affine multiplicative group. Remaining infinitely
many represented terms must occur on strata with positive-dimensional
geometric fibers. This is a restriction on the complete witness
domain, not merely on its generic function field.

\subsection{A smooth birational example with every field value}

\begin{proposition}\label{prop:sharp}
Over any characteristic-zero field $k$, put $E:y^2=x^3+9$ and
let $V\subset\A^4_k$, with coordinates $(t,s,y,z)$, be given by
\begin{equation}\label{eq:surface}
                   y^2=t^3s^3+9,\qquad z^2=s^3+9.
\end{equation}
The surface $V$ is smooth and geometrically irreducible. The map
\[
       \pi:V\longrightarrow E\times E,\qquad
       (t,s,y,z)\longmapsto((ts,y),(s,z))
\]
is birational onto the product of the affine elliptic charts.
Its non-quasi-finite locus is exactly $s=0$, and the regular output
$h=t$ satisfies $h(V(k))=k$.
\end{proposition}
\begin{proof}
On $s\ne0$, the map is an isomorphism onto the open subset where
the second abscissa is nonzero. Its inverse is
\[
             ((x_1,y),(x_2,z))\longmapsto(x_1/x_2,x_2,y,z).
\]
That open subset is smooth and geometrically irreducible. When
$s=0$, equations~\eqref{eq:surface} give $y,z\in\{3,-3\}$,
with $t$ arbitrary. The partial derivatives with respect to $y,z$
form a diagonal submatrix with nonzero entries $2y,2z$, so the
surface is smooth there as well.

Every irreducible component of the two-equation locus in $\A^4$
has dimension at least two. Its part with $s=0$ has dimension one,
so every component meets $s\ne0$. The irreducibility of that open
subset therefore proves geometric irreducibility of $V$.

A fiber with $x_2\ne0$ is a singleton. If $x_2=0$ and $x_1\ne0$
it is empty. Above each of the four pairs with $x_1=x_2=0$ the
fiber is an affine line with coordinate $t$. These are exactly the
non-quasi-finite fibers. In particular $(t,0,3,3)\in V(k)$ for
every $t\in k$, proving surjectivity.
\end{proof}

The equation $x_1-tx_2=0$ is generically a degree-one equation for
$t$, but both coefficients vanish at the four exceptional pairs.
The surface is the affine ratio chart of the blow-up at these
base points. Its exceptional lines explain the failure of an
unguarded ratio formula. They introduce arbitrary field values
even if the elliptic points themselves are restricted to a smaller
field. They also have unbounded denominators at every finite place.

\section{Unirational fibers and bounded denominators}

\begin{lemma}\label{lem:rational-unbounded}
Let $f\in k(T)$ be nonconstant. At infinitely many finite places
of $k$, the valuations of $f(t)$ for $t\in k$ in any fixed
nonempty Zariski-open domain are unbounded below.
Consequently a nonconstant rational function on a $k$-unirational
variety cannot have uniformly bounded denominators on its rational
points in its domain.
\end{lemma}
\begin{proof}
Write $f=F/G$ with coprime polynomials. If $G$ is constant, then
outside finitely many coefficient exceptions the leading term
dominates when $v(t)$ is sufficiently negative, and $v(f(t))$
tends to $-\infty$.

Otherwise choose an irreducible factor $H$ of $G$. After clearing
denominators, its norm to $\Q$ is a nonconstant integer polynomial.
Such a polynomial has prime divisors of its values at infinitely
many rational primes. For completeness, if these primes formed
a finite set $P$, choose an integer $a$ where the polynomial
$Q$ is nonzero and put $D=Q(a)\prod_{p\in P}p$. Then
$Q(a+DT)/Q(a)$ is a nonconstant integer polynomial congruent to
$1$ modulo every prime in $P$. A value of absolute value greater
than one gives a prime divisor outside $P$, a contradiction.

Exclude the finitely many places coming from coefficient
denominators, leading coefficients, and the nonzero resultants
of $H$ with $H'$ and with $F$. Infinitely many remaining prime
divisors of the norm give a place $v$ where $H$ has a simple
residue root. Hensel's lemma gives a root $\alpha\in k_v$ of
$H$, while $F(\alpha)\ne0$. Thus $f$ has a pole at $\alpha$.
Density of $k$ in $k_v$ supplies arguments approaching that pole,
with values of arbitrarily negative valuation. A fixed finite
set of excluded arguments does not alter the conclusion.

For the final assertion, compose a dominant rational
parametrization by projective space with the given function.
The result is nonconstant. Restricting to a suitable $k$-line
gives a nonconstant one-variable rational function, defined away
from finitely many points. Apply the first part.
\end{proof}

\begin{lemma}\label{lem:finite-locus}
If $W\subset A\times\A^1$ is constructible, then the set
\[
             B_0=\{a\in A:W_a\text{ is finite geometrically}\}
\]
is constructible. Moreover $W\cap(B_0\times\A^1)$ is a
constructible set with finite geometric fibers over $A$.
\end{lemma}
\begin{proof}
Decompose $W$ into finitely many locally closed subvarieties.
The dimension of the fibers of each projection has constructible
level sets \cite[Tag 05F9]{Stacks}. A constructible subset of a
geometric affine line is finite exactly when its dimension is at
most zero, including the empty case. Taking finite unions gives
the assertion.
\end{proof}

\begin{theorem}\label{thm:unirational}
Let $\pi:V\to A$ be a morphism to an abelian variety and let
$h:U\to\A^1$ be a rational output on an open subset. Suppose
that every fiber above $a\in A(k)$ meeting $U(k)$ is geometrically
integral and $k$-unirational. If $T=h(U(k))$ has uniformly bounded
denominators, then $T$ has finite intersection with every affine
set $a+b\Lambda$ as in Corollary~\ref{cor:affine-units}.
In particular it cannot contain $\Z$.
\end{theorem}
\begin{proof}
Let $W$ be the constructible image of $(\pi,h)$. For a fiber
contributing a rational witness, the restriction of $h$ is a
rational function on a geometrically integral unirational variety.
If it were nonconstant, Lemma~\ref{lem:rational-unbounded} would
contradict the bound on the denominators of $T$. Thus its entire
geometric image on the domain is a single value.

Every actual rational image point therefore lies in
$W\cap(B_0\times\A^1)$, where $B_0$ is as in
Lemma~\ref{lem:finite-locus}. This is constructible with finite
geometric fibers. Apply Lemma~\ref{lem:constructible}, first to
$h$ and then to each affine transform $(h-a)/b$.
\end{proof}

Smooth conics and positive-dimensional smooth quadrics with a
$k$-point are $k$-rational, by projection from that point.
Theorem~\ref{thm:unirational} therefore applies to families of
such fibers over an abelian base. It also applies to unrestricted
affine fibers. All equations imposed on auxiliary witnesses must
be included in $V$ before the hypothesis is checked.

Finitely many local preparation conditions do not by themselves
remove the denominator problem in Lemma~\ref{lem:rational-unbounded}.
Its proof supplies infinitely many places of unboundedness. Choose
one outside the fixed places, and use weak approximation to approach
its pole while retaining the prescribed nonempty open conditions
elsewhere. A global denominator bound requires a further arithmetic
restriction.

\section{Genus-one families over an abelian base}

We use two established arithmetic facts. Siegel's theorem implies
that for a smooth projective curve $C/k$ of positive genus and a
nonconstant rational function $f$ on $C$, only finitely many points
in the domain have $f$-values in a fixed ring of $S$-integers
\cite[Section 1]{VojtaSurvey}. To see this formulation, use the
finite map $C\to\PP^1$ defined by $f$. After enlarging $S$, its
inverse image of $\A^1$ has a model finite over $\A^1_{\Ok[S^{-1}]}$;
a point with integral $f$-value is integral in that model by
properness of the finite map. Siegel's theorem applies to this
affine positive-genus curve.

Merel's uniform torsion theorem supplies an integer $N=N(k)>0$
that annihilates the rational torsion subgroup of every elliptic
curve over $k$ \cite[Introductory Corollary]{Merel}. One can take
a common multiple of all positive integers below the uniform
bound for the order of a torsion point.

We first describe how a relative divisor replaces the choice of an
origin. For a smooth projective family of geometrically connected
curves, the identity component of the relative Picard scheme is its
relative Jacobian; for genus one it is an elliptic scheme
\cite[Section 9.4, Theorem 1]{BLR}.

\begin{lemma}\label{lem:divisor-map}
Let $\mathcal C\to B$ be a smooth projective family of geometrically
connected genus-one curves, and let $\mathcal L$ be a line bundle
of constant positive degree $d$ on its geometric fibers. Write
$\mathcal J=\operatorname{Pic}^0_{\mathcal C/B}$ for its relative
Jacobian. The formula
\[
 \alpha_{\mathcal L}:\mathcal C\longrightarrow\mathcal J,
 \qquad P\longmapsto
 [\mathcal O_{\mathcal C_b}(dP)\otimes\mathcal L_b^{-1}]
\]
defines a finite morphism whose degree on every geometric fiber is
$d^2$.
\end{lemma}
\begin{proof}
On $\mathcal C\times_B\mathcal C$, the diagonal is a relative
effective Cartier divisor. The line bundle
$\mathcal O(d\Delta)\otimes\operatorname{pr}_2^*\mathcal L^{-1}$
has degree zero on the fibers over the first factor. It therefore
defines the displayed morphism to the relative Picard scheme.

On a geometric fiber choose any origin $O$. Identifying that curve
with its Jacobian, the morphism becomes
\[
             P\longmapsto[d](P-O)
               -[\mathcal L_b\otimes\mathcal O(-dO)].
\]
This is a translate of multiplication by $d$, so it is finite
of degree $d^2$; this divisor-class construction is also described
in \cite[Section 1]{AKMMMP}. The relative morphism is proper,
since $\mathcal C$ is proper over $B$ and $\mathcal J$ is separated
over $B$. Its geometric fibers are finite, so it is finite.
\end{proof}

\begin{theorem}\label{thm:elliptic-family}
Let $B$ be a locally closed subvariety of an abelian variety $A/k$,
and let $\mathcal C\to B$ be a smooth proper family of geometrically
connected genus-one curves, equipped with a relatively ample line
bundle. Let $h:U\to\A^1$ be a rational function on its total space,
on an open domain $U$. If $T=h(U(k))$ has uniformly bounded
denominators, then $T$ has finite intersection with every set
$a+b\Lambda$ of Corollary~\ref{cor:affine-units}. Hence $T$
cannot contain $\Z$. No section of $\mathcal C\to B$ is required.
\end{theorem}
\begin{proof}
The relatively ample bundle makes the proper morphism projective.
Let $W$ be the constructible image of $U$ in $A\times\A^1$.
Separate its finite-fiber locus using
Lemma~\ref{lem:finite-locus}. Outputs above that locus satisfy
Lemma~\ref{lem:constructible}. This includes empty domains on a
fiber. On a nonempty domain in a geometric genus-one fiber, the
image is finite exactly when the induced function is constant.

Consider a remaining fiber over $b\in B(k)$ that contributes a
rational witness. The function induced by $h$ on $\mathcal C_b$
is nonconstant, and its domain omits only finitely many geometric
points. The fixed denominator bound and Siegel's theorem make
$U(k)\cap\mathcal C_b(k)$ finite. This fiber has a rational point,
so choosing one identifies its rational points with those of its
Jacobian $\mathcal J_b$. If $\mathcal J_b(k)$ had positive rank,
the cofinite open part would have infinitely many rational points,
a contradiction. Thus $\mathcal J_b(k)$ consists entirely of
torsion points.

Let $\mathcal L$ be the given relatively ample bundle. Its degree
is a positive locally constant integer; splitting the finitely
many connected components of $B$, assume it equals $d$.
Choose $N=N(k)$ as above and use Lemma~\ref{lem:divisor-map}.
For every rational witness $P$ over a nonconstant contributing fiber,
$\alpha_{\mathcal L}(P)$ belongs to $\mathcal J_b(k)$ and hence
to $\mathcal J_b[N]$. All such witnesses therefore lie in
\[
                 \mathcal F=
                 \alpha_{\mathcal L}^{-1}(\mathcal J[N]).
\]
Both $\alpha_{\mathcal L}$ and $\mathcal J[N]\to B$ are finite,
so $\mathcal F\to B$ is finite. Its composite map to $A$ is
quasi-finite. Apply Theorem~\ref{thm:main} to $h$ on the reduced
scheme underlying $\mathcal F\cap U$. Together with the
finite-output part, this covers all actual rational witnesses.
Apply the same argument to $(h-a)/b$ for the affine assertion;
its denominators remain uniformly bounded after multiplication
by a fixed integer.
\end{proof}

For an elliptic scheme with zero section $O$, take
$\mathcal L=\mathcal O(O)$ and $d=1$; then
$\alpha_{\mathcal L}$ is the identity after the usual identification
with its Jacobian. For a genus-one torsor the divisor-class map
avoids making a varying arithmetic choice of an origin. The finite
subscheme $\mathcal F$ replaces the torsion subscheme of an elliptic
scheme.

The theorem also applies to any open part of such a smooth projective
family, and to finitely many such families on locally closed strata
of the base. Extra witness equations must be included in the source
before applying it. It does not assert that every family with
finitely many rational points on each fiber has a finite geometric
subscheme containing them. In particular the argument does not
extend to higher-genus families merely by applying Faltings's theorem:
their Jacobians may have positive rank even when the curves have
only finitely many rational points.

\section{Consequences for defining rings of integers}

Recall that a subset of $k^m$ is \emph{Diophantine over $k$} if it is
the image under a coordinate projection of the common zero set in
$k^{m+n}$ of finitely many polynomials with coefficients in $k$,
for some $n\ge0$.

For a number field $k$, a Diophantine set satisfying
\eqref{eq:sandwich} would suffice to define $\Ok$ in $k$.
Indeed
\[
                   S'=\bigcup_{j=0}^{M-1}(j+MS)
\]
satisfies $\Z\subseteq S'\subseteq\Ok$. For an integral basis
$\omega_1,\ldots,\omega_d$ one then has
\[
                     \Ok=\omega_1S'+\cdots+\omega_dS'.
\]
Both finite unions and these sums preserve Diophantine definability.

The preceding theorems exclude this strategy when all output values
come from the finite correspondences of Theorem~\ref{thm:main}.
Under the denominator hypothesis they also exclude the unirational
and genus-one families specified above. These obstructions hold over
every number field, independently of its real or complex embeddings.
For a generically finite description, the exceptional loci must be
included in the analysis: they can contain every represented input,
as Proposition~\ref{prop:sharp} demonstrates.

No assertion here concerns arbitrary existential varieties. In
particular, the results do not produce a number field whose ring of
integers is proved not to be Diophantine in the field. They identify
where a construction must introduce additional arithmetic: through
output-varying positive-dimensional fibers that survive the stated
denominator obstructions, or through a different source of witnesses.

\section{Exact algebraic verification}

The standard-library program
\nolinkurl{check_exceptional_correspondence.py} records the exact
identities in Proposition~\ref{prop:sharp}: the two curve equations,
the inverse on $s\ne0$, all four exceptional lines, and the
nonzero boundary Jacobian minors. It also records the symbolic
family $(t,0,3,3)$, whose output is unrestricted. These finite
checks verify the displayed example. The general finiteness and
denominator theorems are established by the proofs and the cited
arithmetic inputs, not by numerical sampling.

\begin{thebibliography}{99}
\bibitem{VojtaI}
P. Vojta, Integral points on subvarieties of semiabelian varieties, I,
\emph{Invent. Math.} \textbf{126} (1996), 133--181,
\url{https://doi.org/10.1007/s002220050092}.

\bibitem{VojtaSurvey}
P. Vojta, Nevanlinna theory and Diophantine approximation,
in \emph{Several Complex Variables}, MSRI Publications \textbf{37},
Cambridge University Press, 1999, 535--564; Theorem 2.2,
\url{https://library.slmath.org/books/Book37/files/vojta.pdf}.

\bibitem{CZ}
P. Corvaja and U. Zannier, with an appendix by D. Masser,
\emph{Rational and integral values of rational functions at rational
points}, arXiv:2608.28255v1, 28 August 2026,
\url{https://arxiv.org/abs/2608.28255v1}.

\bibitem{Merel}
L. Merel, Bornes pour la torsion des courbes elliptiques sur les
corps de nombres, \emph{Invent. Math.} \textbf{124} (1996), 437--449,
\url{https://perso.imj-prg.fr/wp-content/uploads/merel-pub/torsion.pdf}.

\bibitem{BLR}
S. Bosch, W. L\"utkebohmert, and M. Raynaud,
\emph{N\'eron Models}, Ergebnisse der Mathematik und ihrer Grenzgebiete,
3. Folge, vol. 21, Springer, 1990, Chapter 9,
\url{https://doi.org/10.1007/978-3-642-51438-8}.

\bibitem{AKMMMP}
S. Y. An, S. Y. Kim, D. C. Marshall, S. H. Marshall,
W. G. McCallum, and A. R. Perlis,
Jacobians of genus one curves,
\emph{J. Number Theory} \textbf{90} (2001), 304--315,
\url{https://archive.math.arizona.edu/wmc/Research/JacobianFinal.pdf}.

\bibitem{Stacks}
The Stacks Project Authors, \emph{The Stacks Project},
Tags 054K (Chevalley's theorem), 02FW (dimensions of fibers),
and 05F9 (constructible fiber-dimension loci),
\url{https://stacks.math.columbia.edu}.
\end{thebibliography}
\end{document}
